\begin{proposition}[可观测最小多项式的非线性阶塌缩：\texorpdfstring{$q=9,10,11,13$}{q=9,10,11,13} 首次出现固定亏格（可审计）]\label{prop:pom-observable-minpoly-order-collapse-q9-13}
\leanverified{paper\_pom\_observable\_minpoly\_order\_collapse\_q9\_13}
沿用定理 \ref{thm:pom-moment-resonance} 的记号，令 \(P_q(\lambda)\in\ZZ[\lambda]\) 为序列 \(m\mapsto S_q(m)\) 的最小整系数递推特征多项式（亦即可观测最小湮灭多项式；注意本文用 \(\Pi_q\) 表示 Perron 投影，故不采用 \(\Pi_q(\lambda)\) 记号以免混淆）。
则在 \(q\in\{9,10,11,13\}\) 上，\(P_q(\lambda)\) 为不可约整系数多项式，其次数分别为
\[
\deg P_9=7,\qquad \deg P_{10}=9,\qquad \deg P_{11}=9,\qquad \deg P_{13}=11,
\]
并且相对差分模态截断的“名义维数” \(2\lfloor q/2\rfloor+1\) 均出现固定亏格
\[
2\Bigl\lfloor\frac{q}{2}\Bigr\rfloor+1-\deg P_q=2.
\]
特别地，\(q=9\) 是该类“阶塌缩”（Hankel 秩下降）在已核验窗内首次显露的最小 \(q\)。
该固定亏格现象表明：高阶 Rényi 碰撞矩并不必然引入同阶数量的新可观测谱模态；相反，可观测子空间可在离散 \(q\) 上发生代数消隐，从而出现可审计的“复杂度反直觉下降”。
\end{proposition}
\begin{proof}
次数与显式表达式见定理 \ref{thm:pom-moment-resonance}。
不可约性可由模素数不可约证书审计：例如表 \ref{tab:fold_collision_observable_minpoly_galois_certificate_q9_13} 的 \(p_{\mathrm{irr}}\) 列给出素数 \(p\) 使得 \(P_q(\lambda)\bmod p\) 仍不可约；从而 \(P_q\) 在 \(\QQ[\lambda]\) 上不可约，进而在 \(\ZZ[\lambda]\) 上不可约（Gauss 引理）。
阶塌缩等式只是把 \(2\lfloor q/2\rfloor+1\) 与上述次数直接相减。证毕。
\end{proof}

\begin{corollary}[有限状态/有限矩阵实现的维数下界：高阶矩序列强迫状态数增长]\label{cor:pom-finite-state-dimension-lower-bound-from-minpoly}
\leanverified{paper\_pom\_finite\_state\_dimension\_lower\_bound\_from\_minpoly}
沿用命题 \ref{prop:pom-observable-minpoly-order-collapse-q9-13} 的设定。
若存在某个维数 \(d\) 与三元组 \((u,B,v)\)（其中 \(B\in\CC^{d\times d}\)、\(u,v\in\CC^{d}\)）使得
\[
S_q(m)=u^{\mathsf T}B^m v\qquad(m\ge 0),
\]
则序列 \(S_q(\cdot)\) 满足阶数 \(\le d\) 的常系数线性递推，从而必有
\[
d\ \ge\ \deg P_q.
\]
特别地，在 \(q\in\{9,10,11,13\}\) 上分别得到
\[
d_{\min}(S_9)\ge 7,\qquad
d_{\min}(S_{10})\ge 9,\qquad
d_{\min}(S_{11})\ge 9,\qquad
d_{\min}(S_{13})\ge 11.
\]
因此任何以有限状态转导器（状态数与窗口长度 \(m\) 无关）或有限维核实现这些高阶矩序列的机制，都不可避免地满足相应的状态数下界。
\end{corollary}
\begin{proof}
对任意表示 \(S_q(m)=u^{\mathsf T}B^m v\)，由 Cayley--Hamilton 定理，矩阵幂序列 \(B^m\) 满足由 \(B\) 的特征多项式诱导的常系数线性递推；取 \(u^{\mathsf T}(\cdot)v\) 得 \(S_q\) 亦满足某个阶数 \(\le d\) 的递推。
而 \(P_q\) 为最小整系数递推特征多项式，故任何递推阶数都不小于 \(\deg P_q\)，从而 \(d\ge \deg P_q\)。
将命题 \ref{prop:pom-observable-minpoly-order-collapse-q9-13} 的次数代入即得各下界。证毕。
\end{proof}

\begin{proposition}[共振窗 \texorpdfstring{$q=9,\dots,17$}{q=9..17} 的不可约性与模素数证书]\label{prop:pom-resonance-window-minpoly-irreducible-q9-17}
\leanverified{paper\_pom\_resonance\_window\_minpoly\_irreducible\_q9\_17}
对共振窗 \(q\in\{9,10,\dots,17\}\)，由表 \ref{tab:fold_collision_moment_recursions_9_17} 的最小整系数递推系数所诱导的特征多项式 \(P_q(\lambda)\in\ZZ[\lambda]\) 在 \(\QQ[\lambda]\) 上不可约。
更精确地，可取素数证书
\[
(p_9,p_{10},p_{11},p_{12},p_{13},p_{14},p_{15},p_{16},p_{17})
=(11,17,37,29,29,37,17,239,31),
\]
使得对每个 \(q\in\{9,\dots,17\}\)，模 \(p_q\) 约化 \(P_q(\lambda)\bmod p_q\) 在 \(\FF_{p_q}[\lambda]\) 上不可约。
\end{proposition}
\begin{proof}
由表 \ref{tab:fold_collision_moment_recursions_9_17} 的递推
\(
S(m)=\sum_{i=1}^{d_q}c_i S(m-i)
\)
可定义首一多项式
\[
P_q(\lambda):=\lambda^{d_q}-\sum_{i=1}^{d_q}c_i\,\lambda^{d_q-i}\in\ZZ[\lambda].
\]
对每个 \(q\in\{9,\dots,17\}\)，选取上式所列素数 \(p_q\) 并在 \(\FF_{p_q}[\lambda]\) 中检验 \(P_q\bmod p_q\) 不可约，即得 \(P_q\) 在 \(\FF_{p_q}\) 上不可约。
由于 \(P_q\) 为本原首一整数多项式，若其在 \(\QQ[\lambda]\) 可约，则在 \(\ZZ[\lambda]\) 可约，从而在任意素数模约化下亦可约；
因此 \(P_q\bmod p_q\) 的不可约性推出 \(P_q\) 在 \(\QQ[\lambda]\) 上不可约。证毕。
\end{proof}

\begin{corollary}[有理 Rényi 维数的分母下界：由不可约度数强制]\label{cor:pom-renyi-dimension-rational-denominator-degree-lb-q9-17}
沿用命题 \ref{prop:pom-renyi-dimension-spectrum} 的几何 Rényi 维数 \(D_q\)（\(q\ge 2\)），并在共振窗 \(q\in\{9,\dots,17\}\) 上令 \(r_q\) 为特征多项式 \(P_q\) 的 Perron 根（因此 \(r_q\) 在 \(\QQ\) 上的代数次数等于 \(\deg P_q\)，由命题 \ref{prop:pom-resonance-window-minpoly-irreducible-q9-17}）。
若对某个 \(q\in\{9,\dots,17\}\) 有
\[
D_q=\frac{a}{b}\in\QQ,\qquad \gcd(a,b)=1,\qquad b\ge 1,
\]
则必有
\[
b\ \ge\ \left\lceil\frac{\deg(r_q)}{2}\right\rceil
=
\left\lceil\frac{\deg(P_q)}{2}\right\rceil.
\]
特别地，结合表 \ref{tab:fold_collision_moment_recursions_9_17} 的递推阶数
\[
(\deg P_9,\deg P_{10},\deg P_{11},\deg P_{12},\deg P_{13},\deg P_{14},\deg P_{15},\deg P_{16},\deg P_{17})
=(7,9,9,13,11,13,11,13,13),
\]
得到对应的分母下界
\[
b\ge 4,5,5,7,6,7,6,7,7
\qquad
(q=9,\dots,17).
\]
\end{corollary}
\leanverified{paper\_pom\_renyi\_dimension\_rational\_denominator\_degree\_lb\_q9\_17}
\begin{proof}
若 \(D_q=\frac{a}{b}\)，由 \(h_q=q\log 2-\log r_q\) 与 \(D_q=\frac{h_q}{(q-1)\log\varphi}\) 得
\[
r_q
=2^q\,\varphi^{-(q-1)a/b}
\in \QQ\!\left(\varphi^{1/b}\right).
\]
由于 \(\QQ(\varphi)=\QQ(\sqrt5)\) 为二次扩张，而添入 \(b\) 次根 \(\varphi^{1/b}\) 使次数至多增大 \(b\) 倍，故
\[
[\QQ(\varphi^{1/b}):\QQ]\le 2b.
\]
于是 \(\deg(r_q)\le 2b\)，即 \(b\ge \lceil\deg(r_q)/2\rceil\)。
在共振窗内，\(P_q\) 不可约且 \(r_q\) 为其根，故 \(\deg(r_q)=\deg(P_q)\)，代入表 \ref{tab:fold_collision_moment_recursions_9_17} 的次数即得数值下界。证毕。
\end{proof}

\begin{conjecture}[全对称伽罗瓦群与素数谱随机性]\label{conj:pom-resonance-minpoly-galois-sd}
对所有足够大的 \(q\)，可观测最小多项式 \(P_q(\lambda)\) 在 \(\QQ[\lambda]\) 上不可约且其伽罗瓦群满足
\[
\mathrm{Gal}(P_q/\QQ)\cong S_{\deg P_q}.
\]
在该假设下，由 Čebotarev 密度定理可预期：当素数 \(p\) 变化时，\(P_q(\lambda)\bmod p\) 的分解型在 \(S_{\deg P_q}\) 的共轭类上渐近等分布；
相应的模 \(p\) 递推动力学在素数谱意义下呈现接近最大混合的分解指纹。
\end{conjecture}

\begin{proposition}[判别式黑名单与分歧素数集：局部谱分歧仅可能发生在 \texorpdfstring{$p\mid \mathrm{Disc}(P_q)$}{p|Disc(Pq)}（\(q=9,10,11,13\)，可审计）]\label{prop:pom-observable-minpoly-discriminant-ramification-q9-13}
\leanverified{paper\_pom\_observable\_minpoly\_discriminant\_ramification\_q9\_13}
对 \(q\in\{9,10,11,13\}\)，记 \(\Delta_q^{\mathrm{disc}}\) 为 \(\mathrm{Disc}(P_q)\) 的平方自由部分（含符号），并定义伴随二次域
\[
K_q^{\mathrm{disc}}:=\QQ\!\left(\sqrt{\Delta_q^{\mathrm{disc}}}\right).
\]
令
\(
\mathcal{B}_q:=\{p:\ p\ \text{为素数且}\ p\mid \mathrm{Disc}(P_q)\}
\)
为分歧素数集。
则：
\begin{enumerate}
  \item \(\mathrm{Disc}(P_q)\) 的素因子分解由式 \eqref{eq:fold_collision_observable_minpoly_discriminants_q9_13} 给出；因此 \(\Delta_q^{\mathrm{disc}}\) 非平凡，且诱导二次域 \(K_q^{\mathrm{disc}}\)。其 \(\Delta_q^{\mathrm{disc}}\) 与 \(\mathcal{B}_q\) 的显式列表见表 \ref{tab:fold_collision_observable_minpoly_ramified_primes_q9_13}。
  \item 对任意素数 \(p\notin\mathcal{B}_q\)，多项式 \(P_q(\lambda)\bmod p\) 在 \(\FF_p[\lambda]\) 上必为平方自由；因而所有“\(\mathrm{repeat?}\)”型局部谱分歧（重根/Jordan 调制）只能发生在有限黑名单 \(\mathcal{B}_q\) 上。
\end{enumerate}
\end{proposition}
\begin{proof}
第 1 点由脚本 \path{scripts/exp_fold_collision_observable_minpoly_galois_audit_q9_13.py} 对定理 \ref{thm:pom-moment-resonance} 中 \(P_q\) 的判别式做符号计算并整因子分解得到，可复算输出即式 \eqref{eq:fold_collision_observable_minpoly_discriminants_q9_13}。
第 2 点是判别式的标准性质：对 \(\mathrm{char}(\FF_p)\nmid \mathrm{Disc}(P_q)\)（等价于 \(p\notin\mathcal{B}_q\)），有
\(\gcd(P_q,\partial_\lambda P_q)=1\) 于 \(\FF_p[\lambda]\)，故 \(P_q\bmod p\) 无重根，从而平方自由。
结合注 \ref{rem:pom-a2-embedding-local-ramification} 的局部判据，可把“异常素数集”在该窗口内完全操作化为 \(\mathcal{B}_q\)。证毕。
\end{proof}
% Auto-generated by scripts/exp_fold_collision_observable_minpoly_galois_audit_q9_13.py
\begin{equation}\label{eq:fold_collision_observable_minpoly_discriminants_q9_13}
\begin{aligned}
\mathrm{Disc}(P_{9})&=2^{14}\cdot 3^{8}\cdot 5^{4}\cdot 7^{2}\cdot 957380184966212108753,\\
\mathrm{Disc}(P_{10})&=-2^{17}\cdot 5^{2}\cdot 11^{2}\cdot 61\cdot 151\cdot 205423\cdot 189414765073425368820264762991,\\
\mathrm{Disc}(P_{11})&=2^{30}\cdot 3^{9}\cdot 5^{2}\cdot 7^{2}\cdot 11\cdot 13\cdot 139\cdot 8271740008589\cdot 15321736685117\cdot 43780331906709869,\\
\mathrm{Disc}(P_{13})&=2^{66}\cdot 3^{18}\cdot 5^{5}\cdot 7^{6}\cdot 11^{2}\cdot 13\cdot 19\cdot 53\cdot 4637\cdot 2744015218247887\cdot 231576910651001698879\cdot 38806805554433676281923.
\end{aligned}
\end{equation}

% Auto-generated by scripts/exp_fold_collision_observable_minpoly_galois_audit_q9_13.py
\begin{table}[H]
\centering
\scriptsize
\caption{Squarefree discriminant parts $\Delta_q^{\mathrm{disc}}$ and ramified prime sets $\mathcal{B}_q=\{p:\ p\mid \mathrm{Disc}(P_q)\}$ for $q\in\{9,10,11,13\}$.}
\label{tab:fold_collision_observable_minpoly_ramified_primes_q9_13}
\begin{tabular}{r l p{0.62\linewidth}}
\toprule
$q$ & $\Delta_q^{\mathrm{disc}}$ & $\mathcal{B}_q$\\
\midrule
9 & $957380184966212108753$ & $\{2,\allowbreak 3,\allowbreak 5,\allowbreak 7,\allowbreak 957380184966212108753\}$\\
10 & $-2\cdot 61\cdot 151\cdot 205423\cdot 189414765073425368820264762991$ & $\{2,\allowbreak 5,\allowbreak 11,\allowbreak 61,\allowbreak 151,\allowbreak 205423,\allowbreak 189414765073425368820264762991\}$\\
11 & $3\cdot 11\cdot 13\cdot 139\cdot 8271740008589\cdot 15321736685117\cdot 43780331906709869$ & $\{2,\allowbreak 3,\allowbreak 5,\allowbreak 7,\allowbreak 11,\allowbreak 13,\allowbreak 139,\allowbreak 8271740008589,\allowbreak 15321736685117,\allowbreak 43780331906709869\}$\\
13 & $5\cdot 13\cdot 19\cdot 53\cdot 4637\cdot 2744015218247887\cdot 231576910651001698879\cdot 38806805554433676281923$ & $\{2,\allowbreak 3,\allowbreak 5,\allowbreak 7,\allowbreak 11,\allowbreak 13,\allowbreak 19,\allowbreak 53,\allowbreak 4637,\allowbreak 2744015218247887,\allowbreak 231576910651001698879,\allowbreak 38806805554433676281923\}$\\
\bottomrule
\end{tabular}
\end{table}


\begin{proposition}[高阶递推的最大伽罗瓦群：\texorpdfstring{$\mathrm{Gal}(P_q/\QQ)\cong S_d$}{Gal(Pq/Q)=Sd}（\(q=9,10,11,13\)，可审计）]\label{prop:pom-observable-minpoly-galois-sd-q9-13}
\leanverified{paper\_pom\_observable\_minpoly\_galois\_sd\_q9\_13}
设 \(d=\deg P_q\)。
则对 \(q\in\{9,10,11,13\}\) 有
\[
\mathrm{Gal}(P_9/\QQ)\cong S_7,\qquad
\mathrm{Gal}(P_{10}/\QQ)\cong S_9,\qquad
\mathrm{Gal}(P_{11}/\QQ)\cong S_9,\qquad
\mathrm{Gal}(P_{13}/\QQ)\cong S_{11}.
\]
\end{proposition}
\begin{proof}
对每个 \(q\in\{9,10,11,13\}\)，表 \ref{tab:fold_collision_observable_minpoly_galois_certificate_q9_13} 给出三个不分歧素数证书：
\begin{itemize}
  \item 存在素数 \(p_{\mathrm{irr}}\notin\mathcal{B}_q\) 使得 \(P_q\bmod p_{\mathrm{irr}}\) 不可约，于是 Frobenius 含 \(d\)-循环，且 \(P_q\) 在 \(\QQ\) 上不可约（从而 \(G_q:=\mathrm{Gal}(P_q/\QQ)\) 传递）；
  \item 存在素数 \(p_{(d-1,1)}\notin\mathcal{B}_q\) 使得分解型为 \((d-1,1)\)，故 \(G_q\) 含 \((d-1)\)-循环，从而 \(G_q\) 为原始群；
  \item 存在素数 \(p_{\mathrm{J}}\notin\mathcal{B}_q\) 使得分解型包含一条素数长度循环：对 \(d=7\) 含 \(3\)-循环、对 \(d=9\) 含 \(5\)-循环、对 \(d=11\) 含 \(7\)-循环；且该素数长度 \(\le d-3\)。
\end{itemize}
由 Jordan 定理，原始群 \(G_q\) 若含素数长度 \(p\)-循环（\(p\le d-3\)），则 \(A_d\subseteq G_q\)。
另一方面，由命题 \ref{prop:pom-observable-minpoly-discriminant-ramification-q9-13}，\(\mathrm{Disc}(P_q)\) 非平方（\(\Delta_q^{\mathrm{disc}}\neq 1\)），从而 \(G_q\not\subseteq A_d\)。
综上 \(G_q=S_d\)。证毕。
\end{proof}
% Auto-generated by scripts/exp_fold_collision_observable_minpoly_galois_audit_q9_13.py
\begin{table}[H]
\centering
\scriptsize
\caption{Mod-$p$ factorization certificates for $\mathrm{Gal}(P_q/\QQ)\cong S_d$ (unramified primes). Degrees list the factor degrees of $P_q\bmod p$.}
\label{tab:fold_collision_observable_minpoly_galois_certificate_q9_13}
\begin{tabular}{r r r l r l r l}
\toprule
$q$ & $d$ & $p_{\mathrm{irr}}$ & degs & $p_{(d-1,1)}$ & degs & $p_{\mathrm{J}}$ & degs\\
\midrule
9 & 7 & 11 & (7) & 17 & (6,1) & 13 & (3,2,1,1)\\
10 & 9 & 17 & (9) & 109 & (8,1) & 101 & (5,3,1)\\
11 & 9 & 37 & (9) & 17 & (8,1) & 19 & (5,4)\\
13 & 11 & 29 & (11) & 61 & (10,1) & 41 & (7,3,1)\\
\bottomrule
\end{tabular}
\end{table}


\begin{theorem}[共振窗特征多项式在 \texorpdfstring{$q=12,13,14,15$}{q=12,13,14,15} 的满对称伽罗瓦性]\label{thm:pom-resonance-minpoly-galois-sd-q12-15}
\leanverified{paper\_pom\_resonance\_minpoly\_galois\_sd\_q12\_15}
沿用定理 \ref{thm:pom-moment-resonance} 的记号，取 \(q\in\{12,13,14,15\}\)，记 \(d_q:=\deg P_q\)（于是 \(d_{12}=d_{14}=13\)、\(d_{13}=d_{15}=11\)）。
则
\[
\Gal(P_{12}/\QQ)\cong S_{13},\qquad
\Gal(P_{13}/\QQ)\cong S_{11},\qquad
\Gal(P_{14}/\QQ)\cong S_{13},\qquad
\Gal(P_{15}/\QQ)\cong S_{11}.
\]
从而 \(P_{12},P_{13},P_{14},P_{15}\) 均不可由根式求解。
\end{theorem}
\begin{proof}
对 \(q=13\) 的结论已由命题 \ref{prop:pom-observable-minpoly-galois-sd-q9-13} 给出。
以下只需处理 \(q\in\{12,14,15\}\)。
记 \(G_q:=\Gal(P_q/\QQ)\le S_{d_q}\) 为在其 \(d_q\) 个根上的置换表示。

由表 \ref{tab:fold_collision_moment_charpoly_irreducible_certificate_q9_17}，存在不分歧素数 \(p_{\mathrm{irr}}\) 使 \(P_q\bmod p_{\mathrm{irr}}\) 不可约；
于是 \(G_q\) 含一个 \(d_q\)-循环，并特别地在 \(d_q\) 点上是传递的。
另一方面，表 \ref{tab:fold_collision_moment_charpoly_galois_certificate_q12_15} 给出不分歧素数 \(p_{(d_q-1,1)}\) 使得
\(
P_q\bmod p_{(d_q-1,1)}
 \)
具有分解型 \((d_q-1,1)\)；
因此 \(G_q\) 含一个 \((d_q-1)\)-循环（并固定一点）。

为完成证明，只需用下述群论事实。
\par
\noindent\textbf{引理.}
设 \(p\ge 5\) 为素数，\(G\le S_p\) 为传递子群。若 \(G\) 同时包含一个 \(p\)-循环与一个 \((p-1)\)-循环（并固定一点），则 \(G=S_p\)。
\par
\noindent
该引理的标准证明可由交换子生成 \(3\)-循环并推出 \(A_p\subseteq G\)，再利用 \((p-1)\)-循环为奇置换得到 \(G\not\subseteq A_p\)，从而 \(G=S_p\)。

将引理应用于 \((p,q)=(13,12)\)、\((13,14)\) 与 \((11,15)\)，得 \(G_q=S_{d_q}\)。
最后，\(S_p\)（\(p\ge 5\)）不可解，故对应 \(P_q\) 不可由根式求解。证毕。
\end{proof}
% Auto-generated by scripts/exp_fold_collision_moment_charpoly_irreducible_q9_17.py
\begin{table}[H]
\centering
\scriptsize
\caption{Unramified mod-$p$ factorization certificates for $P_q(\lambda)$ with $q\in\{12,13,14,15\}$. The pattern $(d)$ certifies a $d$-cycle (hence transitivity), and $(d-1,1)$ certifies a $(d-1)$-cycle.}
\label{tab:fold_collision_moment_charpoly_galois_certificate_q12_15}
\begin{tabular}{r r r l r l}
\toprule
$q$ & $d_q$ & $p_{\mathrm{irr}}$ & degs mod $p_{\mathrm{irr}}$ & $p_{(d-1,1)}$ & degs mod $p_{(d-1,1)}$\\
\midrule
12 & 13 & 29 & (13) & 17 & (12,1)\\
13 & 11 & 29 & (11) & 61 & (10,1)\\
14 & 13 & 37 & (13) & 47 & (12,1)\\
15 & 11 & 17 & (11) & 127 & (10,1)\\
\bottomrule
\end{tabular}
\end{table}


\begin{theorem}[判别式二次特征的角色独立性与分裂域直积伽罗瓦群]\label{thm:pom-resonance-disc-squareclass-independence-q12-15}
对 \(q\in\{12,13,14,15\}\) 记 \(L_q\) 为 \(P_q\) 的分裂域，并记 \(\Delta_q:=\Disc(P_q)\in\ZZ\)。
则四个平方类
\[
[\Delta_{12}],\ [\Delta_{13}],\ [\Delta_{14}],\ [\Delta_{15}]\in \QQ^\times/(\QQ^\times)^2
\]
在线性空间 \(\QQ^\times/(\QQ^\times)^2\) 的 \(\FF_2\)-结构下线性无关。
因此四个伽罗瓦扩张 \(L_{12},L_{13},L_{14},L_{15}\) 在线性意义下彼此无交，并且
\[
\Gal(L_{12}L_{13}L_{14}L_{15}/\QQ)\cong S_{13}\times S_{11}\times S_{13}\times S_{11}.
\]
\end{theorem}
\begin{proof}
由定理 \ref{thm:pom-resonance-minpoly-galois-sd-q12-15}，对每个 \(q\in\{12,13,14,15\}\) 有 \(\Gal(L_q/\QQ)\cong S_{d_q}\)（\(d_q\in\{13,11\}\)）。
因此 \(L_q\) 的唯一非平凡真正规子群对应于 \(A_{d_q}\triangleleft S_{d_q}\)，从而 \(L_q\) 的唯一二次子域为
\[
K_q:=L_q^{A_{d_q}}=\QQ(\sqrt{\Delta_q}),
\]
其中二次特征由判别式平方类决定。

由表 \ref{tab:fold_collision_moment_charpoly_disc_squareclass_independence_q12_15} 可审计地得到：在素数集 \(\{31,37,43,61\}\) 上，
四个向量
\[
v_q:=\Bigl(\Bigl(\frac{\Delta_q}{31}\Bigr),\Bigl(\frac{\Delta_q}{37}\Bigr),\Bigl(\frac{\Delta_q}{43}\Bigr),\Bigl(\frac{\Delta_q}{61}\Bigr)\Bigr)\in\{\pm 1\}^4
\]
在把 \((-1)\mapsto 1,(+1)\mapsto 0\) 视作 \(\FF_2\) 编码后满秩（秩为 \(4\)）。
由于对固定有限素数集的 Kronecker/Legendre 符号只依赖平方类，这蕴含 \([\Delta_{12}],\dots,[\Delta_{15}]\) 在 \(\QQ^\times/(\QQ^\times)^2\) 中线性无关。

接下来证明分裂域的线性无关性。
令 \(M\) 为其余分裂域的合成域，则 \(L_q\cap M\) 作为 \(L_q\) 的正规子扩张，必为 \(\QQ\) 或 \(K_q\)。
若 \(K_q\subseteq M\)，则 \([\Delta_q]\) 属于由其余 \([\Delta_{q'}]\) 生成的子空间，违背平方类线性无关。
故 \(L_q\cap M=\QQ\)。
对 \(q=12,13,14,15\) 逐步递推得到四个分裂域两两及整体线性无关。

最后，对伽罗瓦扩张的线性无关性应用标准合成律，得到合成分裂域的伽罗瓦群为直积同构。证毕。
\end{proof}
% Auto-generated by scripts/exp_fold_collision_moment_charpoly_irreducible_q9_17.py
\begin{table}[H]
\centering
\scriptsize
\setlength{\tabcolsep}{6pt}
\caption{Discriminant squareclass fingerprints for $q\in\{12,13,14,15\}$ using Legendre symbols at primes $p\in\{31,37,43,61\}$. The induced $4\times 4$ binary matrix (encoding $-1\mapsto 1$ and $+1\mapsto 0$) has rank 4.}
\label{tab:fold_collision_moment_charpoly_disc_squareclass_independence_q12_15}
\begin{tabular}{r r r r r}
\toprule
$q$ & $\left(\frac{\mathrm{Disc}(P_q)}{31}\right)$ & $\left(\frac{\mathrm{Disc}(P_q)}{37}\right)$ & $\left(\frac{\mathrm{Disc}(P_q)}{43}\right)$ & $\left(\frac{\mathrm{Disc}(P_q)}{61}\right)$\\
\midrule
12 & -1 & -1 & 1 & 1\\
13 & 1 & -1 & 1 & -1\\
14 & -1 & 1 & 1 & 1\\
15 & -1 & -1 & -1 & 1\\
\bottomrule
\end{tabular}
\end{table}


\begin{corollary}[两条碰撞矩谱域的线性互斥与直积伽罗瓦群（\texorpdfstring{$q=14,15$}{q=14,15}）]\label{cor:pom-resonance-linear-disjointness-q14-15}
\leanverified{paper\_pom\_resonance\_linear\_disjointness\_q14\_15}
设 \(L_{14},L_{15}\) 分别为 \(P_{14},P_{15}\) 的分裂域。
则
\[
L_{14}\cap L_{15}=\QQ,
\qquad
\Gal(L_{14}L_{15}/\QQ)\cong S_{13}\times S_{11}.
\]
\end{corollary}
\begin{proof}
由定理 \ref{thm:pom-resonance-disc-squareclass-independence-q12-15} 的证明，对任意 \(q\in\{12,13,14,15\}\)，令 \(M\) 为其余分裂域的合成域则 \(L_q\cap M=\QQ\)。
取 \(q=14\) 且 \(M=L_{12}L_{13}L_{15}\) 得 \(L_{14}\cap L_{15}=\QQ\)。
又由定理 \ref{thm:pom-resonance-minpoly-galois-sd-q12-15}，\(L_{14}/\QQ\) 与 \(L_{15}/\QQ\) 为伽罗瓦扩张，且其伽罗瓦群分别同构于 \(S_{13}\) 与 \(S_{11}\)；
在交域为 \(\QQ\) 时，标准合成律给出 \(\Gal(L_{14}L_{15}/\QQ)\cong \Gal(L_{14}/\QQ)\times\Gal(L_{15}/\QQ)\)。证毕。
\end{proof}

\begin{corollary}[四多项式分解型的切博塔廖夫完全独立律与显式密度]\label{cor:pom-resonance-chebotarev-product-density-q12-15}
在定理 \ref{thm:pom-resonance-disc-squareclass-independence-q12-15} 的设定下，
对不分歧素数 \(p\) 记 \(\mathrm{Frob}_p\) 为 \(\Gal(L_{12}L_{13}L_{14}L_{15}/\QQ)\) 的 Frobenius 共轭类。
则对任意共轭类 \(C_q\subseteq S_{d_q}\)（\(q\in\{12,13,14,15\}\)），满足
\[
\mathrm{Frob}_p\in C_{12}\times C_{13}\times C_{14}\times C_{15}
\]
之素数集合具有自然密度
\[
\delta=\prod_{q\in\{12,13,14,15\}}\frac{|C_q|}{|S_{d_q}|}.
\]
特别地，四个事件 ``\(P_q\bmod p\) 在 \(\FF_p[\lambda]\) 中不可约'' 彼此独立，
其同时发生的密度为
\[
\frac1{13}\cdot\frac1{11}\cdot\frac1{13}\cdot\frac1{11}=\frac1{20449}.
\]
\end{corollary}
\begin{proof}
由定理 \ref{thm:pom-resonance-disc-squareclass-independence-q12-15}，合成分裂域的伽罗瓦群为直积。
对该伽罗瓦扩张应用切博塔廖夫密度定理，即得对任意直积共轭类的密度公式。
再取 \(C_q\) 为 \(S_{d_q}\) 中 \(d_q\)-循环的共轭类（其大小为 \((d_q-1)!\)），得到不可约事件密度为 \(1/d_q\)，并由直积结构得到乘积分解。证毕。
\end{proof}

\begin{theorem}[四个 Perron 谱半径的强乘法独立性]\label{thm:pom-resonance-perron-multiplicative-independence-q12-15}
对 \(q\in\{12,13,14,15\}\) 记 \(r_q\) 为与 \(P_q\) 同特征多项式的共振窗算子的 Perron 根，
并令
\[
\Lambda_q:=\max\{|\mu|:\ P_q(\mu)=0,\ \mu\neq r_q\}.
\]
则在表 \ref{tab:fold_collision_kernel_rh_scan_q2_17} 的已核验范围内，对这些 \(q\) 都有 \(\Lambda_q<r_q\)。
在该严格谱隙条件下，若
\[
r_{12}^{a}\,r_{13}^{b}\,r_{14}^{c}\,r_{15}^{d}\in\QQ^\times
\qquad(a,b,c,d\in\ZZ),
\]
则必有 \(a=b=c=d=0\)。
\end{theorem}
\begin{proof}
设存在 \((a,b,c,d)\neq (0,0,0,0)\) 使 \(r_{12}^{a}r_{13}^{b}r_{14}^{c}r_{15}^{d}\in\QQ^\times\)，取其中指数非零的某一 \(q_0\in\{12,13,14,15\}\)。
由定理 \ref{thm:pom-resonance-minpoly-galois-sd-q12-15} 与定理 \ref{thm:pom-resonance-disc-squareclass-independence-q12-15}，
可取 \(\sigma\in\Gal(L_{12}L_{13}L_{14}L_{15}/\QQ)\) 在除 \(L_{q_0}\) 以外的三因子上恒等，
而在 \(L_{q_0}\) 上把 \(r_{q_0}\) 送到某一不同共轭根 \(r'_{q_0}\)。
由于 \(r_{q_0}\) 为唯一最大模长根且 \(\Lambda_{q_0}<r_{q_0}\)（表 \ref{tab:fold_collision_kernel_rh_scan_q2_17}），必有 \(|r'_{q_0}|\le \Lambda_{q_0}<r_{q_0}\)。

对关系式施加 \(\sigma\) 并利用其固定 \(\QQ^\times\) 与其余三项，得到
\(
(r'_{q_0})^{n}=r_{q_0}^{n}
\)
对某个 \(n\in\ZZ\setminus\{0\}\) 成立，取模长即得
\(
|r'_{q_0}|^{|n|}=r_{q_0}^{|n|},
\)
与 \(|r'_{q_0}|<r_{q_0}\) 矛盾。证毕。
\end{proof}

\begin{conjecture}[共振窗分裂域的判别式平方类无限独立性]\label{conj:pom-resonance-disc-squareclass-infinite-independence}
对共振窗族 \(\{P_q\}_{q\ge 9}\)，设对所有足够大的 \(q\) 有 \(\Gal(P_q/\QQ)\cong S_{\deg P_q}\)，并且任意有限子集 \(\{q_1,\dots,q_m\}\) 的判别式平方类
\(
[\Disc(P_{q_i})]\in \QQ^\times/(\QQ^\times)^2
\)
线性无关。
则相应分裂域在 \(\QQ\) 上两两及整体线性无关，从而任意有限合成分裂域的伽罗瓦群为直积 \(\prod_i S_{\deg P_{q_i}}\)，并导出对多 \(q\) 同时分解型的切博塔廖夫完全独立律。
\end{conjecture}

\begin{corollary}[Chebotarev 密度的可检验频率律：不可约密度为 \texorpdfstring{$1/d$}{1/d}]\label{cor:pom-observable-minpoly-chebotarev-density-q9-13}
在命题 \ref{prop:pom-observable-minpoly-galois-sd-q9-13} 的结论下，对去除有限分歧素数集 \(\mathcal{B}_q\) 后的素数集合，任意分裂型（等价于 \(S_d\) 中某个共轭类）出现的自然密度等于该共轭类在 \(S_d\) 中的比例。
特别地，\(P_q(\lambda)\bmod p\) 仍不可约（Frobenius 为 \(d\)-循环）的密度为 \(1/d\)，故
\[
\mathrm{dens}\{p:\ P_9\bmod p\ \text{不可约}\}=\frac17,\quad
\mathrm{dens}\{p:\ P_{10}\bmod p\ \text{不可约}\}=\frac19,
\]
\[
\mathrm{dens}\{p:\ P_{11}\bmod p\ \text{不可约}\}=\frac19,\quad
\mathrm{dens}\{p:\ P_{13}\bmod p\ \text{不可约}\}=\frac1{11}.
\]
\end{corollary}
\begin{proof}
由 Chebotarev 密度定理，未分歧素数的 Frobenius 共轭类在 \(G_q=S_d\) 中按共轭类大小等分布；而 \(d\)-循环在 \(S_d\) 中所占比例为 \((d-1)!/d!=1/d\)。证毕。
\end{proof}
\leanverified{paper\_pom\_observable\_minpoly\_chebotarev\_density\_q9\_13}

\begin{corollary}[判别式二次特征给出 Frobenius 奇偶的闭式判据]\label{cor:pom-observable-minpoly-frob-parity-q9-13}
\leanverified{paper\_pom\_observable\_minpoly\_frob\_parity\_q9\_13}
在 \(G_q=S_d\) 下，\(S_d\) 的唯一指数 \(2\) 正规子群为 \(A_d\)，其固定子域为 \(\QQ(\sqrt{\mathrm{Disc}(P_q)})=\QQ(\sqrt{\Delta_q^{\mathrm{disc}}})\)。
因此对任意不分歧素数 \(p\notin\mathcal{B}_q\)，Frobenius 置换的奇偶性由 Kronecker 符号严格给出：
\[
\mathrm{sgn}(\mathrm{Frob}_p)=\left(\frac{\Delta_q^{\mathrm{disc}}}{p}\right).
\]
特别地，对 \(q=9\)，由式 \eqref{eq:fold_collision_observable_minpoly_discriminants_q9_13} 可知 \(\Delta_9^{\mathrm{disc}}\) 为素数，从而奇偶判据退化为单一 Legendre 符号 \(\left(\frac{\Delta_9^{\mathrm{disc}}}{p}\right)\)。
\end{corollary}
\begin{proof}
这是判别式的标准刻画：\(\mathrm{Disc}(P_q)\) 的平方类对应 \(G_q\) 到 \(\{\pm 1\}\) 的符号表示（置换奇偶）所切出的二次子域；对不分歧素数，Frobenius 的像由二次互反律的 Kronecker 符号读出。证毕。
\end{proof}

\begin{proposition}[负实主导律与交替修正：\texorpdfstring{$q=9,10,11,13$}{q=9,10,11,13} 的有效 Ramanujan 比率]\label{prop:pom-observable-minpoly-negative-real-dominance-q9-13}
对 \(q\in\{9,10,11,13\}\)，记 \(r_q\) 为 \(P_q\) 的 Perron 根，并令
\[
\Lambda_q:=\max\{|\mu|:\ P_q(\mu)=0,\ \mu\neq r_q\}.
\]
则达到 \(\Lambda_q\) 的根为唯一负实根 \(\mu^-_q<0\)，从而 \(S_q(m)\) 的首个非主指数修正带严格交替相位 \(({-1})^m\)。
并且有效 Ramanujan 比率
\(
R_q:=|\mu^-_q|/\sqrt{r_q}
\)
在该窗口内满足式 \eqref{eq:fold_collision_observable_minpoly_negative_real_dominance_q9_13} 的可复算数值（均严格大于 \(1\)）。
\end{proposition}
\begin{proof}
由脚本 \path{scripts/exp_fold_collision_observable_minpoly_galois_audit_q9_13.py} 对 \(P_q\) 的全部根做高精度数值求解，并逐一比较模长得到：除 Perron 根外，最大模根为唯一负实根。
相应地，递推解的谱展开中该根贡献的符号为 \((\mu^-_q)^m=(-1)^m|\mu^-_q|^m\)，故出现交替修正；比率 \(R_q\) 的数值即由该根与 \(r_q\) 直接代入定义得到。证毕。
\end{proof}
% Auto-generated by scripts/exp_fold_collision_observable_minpoly_galois_audit_q9_13.py
\begin{equation}\label{eq:fold_collision_observable_minpoly_negative_real_dominance_q9_13}
\begin{aligned}
R_{9}:=\frac{|\mu^-_{9}|}{\sqrt{r_{9}}}&\approx 2.05868056993\qquad(r_{9}\approx 11.778421934988905909,\ \mu^-_{9}\approx -7.0653311115125455995),\\
R_{10}:=\frac{|\mu^-_{10}|}{\sqrt{r_{10}}}&\approx 2.3677822786\qquad(r_{10}\approx 14.771098354824257587,\ \mu^-_{10}\approx -9.1001418241828733831),\\
R_{11}:=\frac{|\mu^-_{11}|}{\sqrt{r_{11}}}&\approx 2.72056851827\qquad(r_{11}\approx 18.535847481911089603,\ \mu^-_{11}\approx -11.712939278722332302),\\
R_{13}:=\frac{|\mu^-_{13}|}{\sqrt{r_{13}}}&\approx 3.58228432437\qquad(r_{13}\approx 29.237338790942123832,\ \mu^-_{13}\approx -19.369971099960025432).
\end{aligned}
\end{equation}

